% 06_conclusion.tex
% Chapter 6: Conclusion

\section{Conclusion}

This project developed and evaluated a Ku-band GEO bent-pipe satellite link from Rome to Ankara via the Eutelsat Hotbird 13G satellite. The core analysis followed the standard link-budget procedure and was implemented in a custom Python framework. Additional modules were then used to examine interference, atmospheric attenuation, dynamic receiver noise, Monte Carlo rain states, and DVB-S2 ACM.

In the baseline clear-sky scenario without interference, the link calculations resulted in a total $(C/N_0)_{\text{total}}$ of \cnTotal~dB-Hz and an $E_b/N_0$ of \ebnoTotal~dB. This provided a \marginClear~dB margin over the 7.0 dB required for the 10 Mbps baseline. However, introducing Adjacent Satellite Interference (ASI) and Intermodulation Distortion (IMD) reduced this margin to \marginASI~dB, showing the importance of accounting for a crowded GEO arc.

The atmospheric propagation analysis, modeled after ITU-R guidelines, showed the significant impact of rain fade at Ku-band frequencies. At a 99.9\% availability target, the combined effects of signal attenuation and the resulting increase in dynamic system noise temperature consumed the remaining margin, leaving the fixed-rate 10 Mbps link with a deficit of \marginRain~dB.

To address this deficit without increasing antenna sizes or transmitter power, a DVB-S2 Adaptive Coding and Modulation (ACM) approach was examined. While the fixed-rate link fails during a 99.9\% rain event, the ACM model demonstrated that the system could step down to a more robust MODCOD (such as QPSK 2/5), preserving the connection at a reduced data rate.

In summary, this project illustrated the step-by-step process of sizing a satellite link and the necessity of looking beyond clear-sky thermal noise. By incorporating basic interference and dynamic weather effects into the Python model, the analysis highlighted how practical satellite networks manage physical constraints to maintain operation.
